\documentclass[11pt]{article}
\usepackage{amsmath, amssymb}
\usepackage{hyperref}
\usepackage{graphicx}
\usepackage{parskip}

% Define \problemname and \illustration if not using Kattis template
\newcommand{\problemname}[1]{\section*{#1}}
\newcommand{\illustration}[3]{%
    \begin{figure}[h]
        \centering
        \includegraphics[width=#1\textwidth]{#2}
        \caption{#3}
    \end{figure}
}

\begin{document}

\problemname{I Can't See the Screen!}

A cinema has $R$ rows and $C$ columns of seats, with row $0$ being farthest from the screen, while row $R-1$ is closest to the screen, at the bottom. The floor is sloped downwards to the screen such that each row is exactly $1$ foot higher than the row in front of it.

You are given an $R \times C$ matrix $H$, where $H[r][c]$ is the height in feet of the person sitting in row $r$, column $c$, or $-1$ if the seat is empty.

\illustration{0.5}{cinema.png}{Photo by Felix Mooneeram from unsplash.com}

\section*{Visibility Rule}

A person sitting in $H[r][c]$ cannot see the screen if there exists a person in $H[k][c]$ in front of them (i.e., $k > r$, closer to the screen) and whose head is higher in absolute position or at the same level, accounting for both their height and the elevation of their row.

Among all such people $k$ who block the view, only the nearest one (the smallest $k > r$) matters.

If no such person exists, the person can see the screen.

\section*{Input}

You are given a line containing two space-separated integers $n, m \le 10\,000$ that represent the dimensions of $H$. Each of the next $n$ lines contains $m$ space-separated integers that represent the height of the person sitting in $H[r][c]$ where $0 < H[r][c] < 10^9$. If the seat is empty, the value is $-1$.

\section*{Output}

Print an $R \times C$ grid of space-separated integers. For each $H[r][c]$, output $-1$ if the seat is empty or if the person can see the screen; otherwise output the row index $k$ of the nearest person in front of them who blocks their view.

\section*{Sample Input 1}

\begin{verbatim}
3 3
4 8 3
4 2 3
3 5 3
\end{verbatim}

\section*{Sample Output 1}

\begin{verbatim}
-1 -1 -1
-1 2 -1
-1 -1 -1
\end{verbatim}

\section*{Sample Input 2}

\begin{verbatim}
3 3
5 1 3
-1 2 3
7 5 -1
\end{verbatim}

\section*{Sample Output 2}

\begin{verbatim}
2 1 -1
-1 2 -1
-1 -1 -1
\end{verbatim}

\end{document}
